Linear probes can distinguish human- from machine-written text in language-model activations, but a readable feature need not be a causal control variable. We test whether a dense residual-stream direction fitted to content-paired human and machine continuations also changes how artificial a model's own writing appears. We estimate authorship directions in Qwen2.5-7B-Instruct and its pretrained base checkpoint using HAP-E, intervene during autoregressive decoding, and score held-out generations with a MAGE-trained lexical detector and an independently trained public RoBERTa detector. Authorship is almost perfectly readable: the validation-selected layer reaches 0.9997 held-out AUROC in both checkpoints. The write-side result is different. Across 16 instruct-model prompts, moving from strength $-2$ to $+2$ changes the MAGE detector's $P(\mathrm{AI})$ by $+0.008$ (95\% cluster-bootstrap CI $[-0.077,0.095]$) and the public detector's score by $-0.077$ ($[-0.246,0.057]$). Neither effect is significant, the detectors disagree in sign, and nuisance-orthogonal, random, shuffled-label, formality, and Assistant-axis controls reveal no specific effect. A 12-prompt base-model replication is likewise null. Word count, repetition, semantic similarity, and perplexity do not indicate endpoint quality collapse. These results establish a read/write dissociation within the tested design: authorship information is linearly accessible, but the validation-selected direction is not a reliable generative control for AI-sounding prose at the tested layer and doses.

